% =====================================================================
% CONFIGURATION DE LA NUMÉROTATION DES TITRES
% =====================================================================


% --- 1. STYLE DES COMPTEURS ---
% Options possibles : \arabic (1, 2...), \Roman (I, II...), \roman (i, ii...), \Alph (A, B...), \alph (a, b...)
\renewcommand{\thepart}{\Alph{section}}
\renewcommand{\thesection}{\arabic{section}.}
\renewcommand{\thesubsection}{\arabic{subsection}\Alph{section}.}
\renewcommand{\thesubsubsection}{\alph{subsubsection})}

% --- 2. PRÉFIXES PERSONNALISÉS ---
% Mots qui précèdent les numéros (laissez les accolades vides {} pour ne rien mettre)
\newcommand{\prefixPart}{Partie }      % Donne "Partie I"
\newcommand{\prefixSection}{}    % Donne "Étape 1"
\newcommand{\prefixSubsection}{}       % Pas de préfixe
\newcommand{\prefixSubsubsection}{}    % Pas de préfixe

% --- 3. SÉPARATEURS ---
% Espace ou ponctuation entre la numérotation et le nom du titre
\newcommand{\sepPart}{~--~}            % Donne "Partie I -- Mon Titre"
\newcommand{\sepSection}{~~}        % Donne "Étape 1    Mon Titre"
\newcommand{\sepSubsection}{\enspace}
\newcommand{\sepSubsubsection}{~}



